\section{Related Work}
\label{sec:related}

\paragraph{Adaptive retrieval.}
Adaptive-RAG routes questions to retrieval strategies using a learned
complexity classifier~\citep{adaptiverag}, while SKR estimates whether a model
already knows an answer before retrieval~\citep{skr}. FLARE instead drafts the
next sentence and retrieves when low-confidence tokens appear~\citep{flare};
DRAGIN combines a generation-time trigger with attention-based query
formulation~\citep{dragin}. TARG studies a lightweight question-level decision
from prefix logits, including entropy and top-logit margins~\citep{targ}.
These methods make uncertainty useful operationally, but most apply it either
before reasoning or at the free-form token/sentence level. \method{} applies a
training-free signal at explicit sub-question boundaries.

Learned policies provide finer control. Self-RAG trains reflection tokens that
control retrieval and evaluate evidence~\citep{selfrag}. DeepRAG explores
retrieve-or-answer decisions at individual sub-queries and calibrates the
policy from preference data~\citep{deeprag}. Such policies can jointly learn
reasoning and retrieval behavior, whereas our objective is complementary: a
score that can be inspected and attached to an existing greedy generation
without policy training.

\paragraph{Decomposition and multi-step retrieval.}
IRCoT interleaves chain-of-thought sentences with retrieval~\citep{ircot}.
Question-decomposition methods instead expose explicit intermediate
operations or sub-questions~\citep{decomprc,break_qdmr}. Decomposition can
create sharper retrieval queries, but a wrong bridge answer can corrupt later
queries and short sub-questions can omit constraints from the original
question. \method{} therefore retains a direct candidate and keeps
original-question evidence available to the decomposition path. We treat these
components as an evaluation framework for the confidence signal rather than
attributing all system gains to decomposition alone.

\paragraph{Confidence and uncertainty in language models.}
Sequence probability, token entropy, verbalized self-evaluation such as
P(True), and disagreement across sampled answers capture different aspects of
model confidence~\citep{kadavath2022,semanticuncertainty,semanticentropy}.
They are useful ranking signals, but a model may remain confidently wrong.
\tailleak{} is closest to logit-based uncertainty methods, yet differs in what
it retains: it ignores the exact probability arrangement within the head and
tail and records only the mass crossing their boundary. This choice follows
from an entropy decomposition in Section~\ref{sec:tail-leakage} and is tested
against representative alternatives in Section~\ref{sec:confidence-results}.
